
\documentclass[11pt]{article}

% ---------- Page setup ----------
\usepackage[a4paper,margin=1in]{geometry}
\usepackage[dvipsnames]{xcolor}
\usepackage{microtype}
\usepackage{parskip} % no paragraph indent, adds vertical space

% ---------- Math packages ----------
\usepackage{amsmath,amssymb,amsthm,mathtools}
\usepackage{bm}       % bold math symbols

% ---------- Theorem environments ----------
\theoremstyle{plain}
\newtheorem{theorem}{Theorem}[section]
\newtheorem{lemma}[theorem]{Lemma}
\newtheorem{proposition}[theorem]{Proposition}
\newtheorem{corollary}[theorem]{Corollary}

\theoremstyle{definition}
\newtheorem{definition}[theorem]{Definition}
\newtheorem{example}[theorem]{Example}

\theoremstyle{remark}
\newtheorem{remark}[theorem]{Remark}
\newtheorem{note}[theorem]{Note}

% ---------- Lists / links ----------
\usepackage{enumitem}
\usepackage[hidelinks]{hyperref}

% ---------- Handy macros (edit to taste) ----------
\newcommand{\R}{\mathbb{R}}
\newcommand{\N}{\mathbb{N}}
\newcommand{\Z}{\mathbb{Z}}
\newcommand{\Q}{\mathbb{Q}}
\newcommand{\C}{\mathbb{C}}
\newcommand{\eps}{\varepsilon}
\newcommand{\dd}{\,\mathrm{d}}

% ---------- Title info ----------
\title{Random Variable}
\author{Joe and Tony}
\date{\today}

\begin{document}
	\maketitle
	\tableofcontents

\newpage
	
	\section{Random variable}
	
	\textcolor{red}{In this note, we denote a probability space as $\Omega$, a $\sigma$-algebra as $\mathcal{F}$, and a probability measure as $\mathbb{P}$. This follows the convention of Prof. Bao.}
	
	Random variable is a numerical quantity associated to physical outcome.
	
	\begin{definition}
		A random variable is a `nice function' from \(\Omega\) to \(\R\),
		\(X:\Omega\to \R,\ \omega\mapsto X(\omega)\).
	\end{definition}
	
	\begin{remark}
		\textcolor{red}{Prof. Bao mentioned in class that it is of theoretical interest to define `nice' function. However, this is not used in the course.}
	\end{remark}
	
	\begin{remark}
		Due to the uncertainty of \(\omega\), the value \(X(\omega)\) is also uncertain. \textcolor{red}{A better way to say this is that $X$ sends the given probability space $(\Omega,\mathcal{F},\mathbb{P})$ on $\Omega$, to some new probability space $(\mathbb{R},X_*(\mathcal{F}),X_*(\mathbb{P}))$ on $\mathbb{R}$. $X_*(\mathcal{F})=\{\Xi\subseteq\mathbb{R}:X^{-1}(\Xi)\in\mathcal{F}\}$. $X_*(\mathbb{P})(\Xi)=\mathbb{P}(X^{-1}(\Xi))$. Prof. Bao mentioned in class that this is called the induced probability space on $\mathbb{R}$.}
	\end{remark}
	
	
	A natural question is: probability of \(X(\omega)=x\) for any \(x\in\R\)?
	
	\begin{example}
		Two fair coins:
		\(\Omega=\{HH,HT,TH,TT\}\).
		The random variable \(X\) of the number of heads is described by the following table.
		$$
		\begin{tabular}
			{c|cc}
			& $H$ & $T$ \\
			\hline
			$H$ & $2$ & $1$ \\
			$T$ & $1$ & $0$ \\
		\end{tabular}
		$$
	\end{example}
	
	\begin{example}
		Toss 2 fair dice:
		\(\Omega=\{(1,1),(1,2),\dots,(6,6)\}\), \(\mathbb{P}(\{\omega\})=\frac{1}{36}\).
		Take \(X\) as sum of 2 dice, \(X\in\{2,3,\dots,12\}\).
		\[
		\mathbb{P}(X=2)=\mathbb{P}(\{\omega: X(\omega)=2\})
		=\mathbb{P}(\{(1,1)\})=\frac{1}{36}.
		\]
	\end{example}
	
	\begin{example}
		Cutting a rope of length \(1\) at a random point:
		\(\Omega=[0,1]\).
		\(L:\Omega\to \R,\ \omega\mapsto L(\omega)\) = length of the shorter part,
		\[
		L(\omega)=\min\{\omega,1-\omega\}.
		\]
	\end{example}
	
	\begin{example}
		Cutting a rope at a random point \(\omega\).
		Let \(X(\omega)=\omega\).
		\[
		\mathbb{P}(X=x)=\mathbb{P}(\{\omega: X(\omega)=x\})=\mathbb{P}(\{x\})=0.
		\]
	\end{example}
	
	\begin{remark}
		In continuous case, \(\mathbb{P}(X=x)\) is not informative to tell the probabilistic behavior of \(X\).
	\end{remark}
	
	Correct question to ask: probability of \(X\) being in any interval?
	\(\mathbb{P}(X\in(a,b))\) or \(\mathbb{P}(X\in[a,b])\) or \(\mathbb{P}(X\in(-\infty,b])\), etc.
	
	\section{Cumulative distribution function (cdf)}
	
	
	
	\begin{definition}
		[cumulative distribution function] Cdf \(F(x)\) of a r.v. \(X\) is defined as
		\[
		F(x)=\mathbb{P}(X\le x)=\mathbb{P}(X\in(-\infty,x])
		=\mathbb{P}(\{\omega\in \Omega:\ X(\omega)\in(-\infty,x]\}).
		\]
	\end{definition}
	
	\begin{remark}
		For any \(a<b\),
		\[
		\mathbb{P}(a<X\le b)=\mathbb{P}(\{\omega:\ a<X(\omega)\le b\})
		=\mathbb{P}(X\le b)-\mathbb{P}(X\le a)=F(b)-F(a).
		\]
		Also,
		\[
		\mathbb{P}(X\in(a,b])=F(b)-F(a).
		\]
	\end{remark}
	
	\begin{remark}
		In the discrete case, if \(X\) takes values on \(\{x_1,x_2,\dots\}\), then
		\[
		F(x)=\mathbb{P}(X\le x)
		=\sum_{x_i\le x}\mathbb{P}(X=x_i).
		\]
	\end{remark}
	
	\begin{proposition}
		\begin{enumerate}
			\item [(i)] For any \(a\le b\), \(F(a)\le F(b)\) (non-decreasing).
			\item [(ii)] \(\displaystyle \lim_{x\to +\infty}F(x)=1\).
			\item [(iii)] \(\displaystyle \lim_{x\to -\infty}F(x)=0\).
			\item [(iv)] For all \(a<b\),
			\[
			\mathbb{P}(a<X\le b)=F(b)-F(a).
			\]
		\end{enumerate}
	\end{proposition}
	
	Q: How to get \(\mathbb{P}(X=a)\) from \(F(x)\)?
	
	\begin{example}
		\(\Omega=\{HH,HT,TH,TT\}\),
		\(X(HH)=2\), \(X(HT)=X(TH)=1\), \(X(TT)=0\).
		\[
		\mathbb{P}(X=2)=\mathbb{P}(\{\omega:\ X(\omega)=2\})=\mathbb{P}(\{HH\})=\frac14.
		\]
	\end{example}
	
	\begin{remark}
		\(F(x)\) is right continuous, i.e.
		\[
		F(x)=\lim_{\delta\downarrow 0}F(x+\delta).
		\]
		It may not be left continuous; but the left limit exists for all \(x\).
		Actually,
		\[
		\lim_{\delta\downarrow 0}F(x-\delta)=\mathbb{P}(X<x)=:F(x-).
		\]
		Hence,
		\[
		\mathbb{P}(X=x)=\mathbb{P}(X\le x)-\mathbb{P}(X<x)=F(x)-F(x-).
		\]
	\end{remark}
	
	\begin{remark}
		\[
		\mathbb{P}(a\le X\le b)
		=\mathbb{P}(a<X\le b)+\mathbb{P}(X=a)
		=F(b)-F(a)+\big(F(a)-F(a-)\big)
		=F(b)-F(a-).
		\]
	\end{remark}
	
	\section{More specific discussions on discrete r.v.}
	
	\begin{definition}
		Probability mass function (p.m.f.) \(f(a)\) of \(X\) is defined by
		\[
		f(a)=\mathbb{P}(X=a),
		\]
		where \(X\) takes values on \(\{a_1,a_2,\dots\}\).
	\end{definition}
	
	\begin{remark}
		\[
		F(b)=\sum_{a_i\le b} f(a_i).
		\]
		Also,
		\[
		f(b)=\mathbb{P}(X=b)=F(b)-F(b-)=F(b)-\lim_{\delta\downarrow 0}F(b-\delta).
		\]
	\end{remark}
	
	Important examples of discrete r.v.
	
	Bernoulli trials: independently repeated experiment with one specific event $\omega$ \textcolor{red}{(should assume that the singleton $\{\omega\}$ lies in the sigma algebra?)} identified as success
	(e.g.1: coin flipping, may not be fair. e.g.2: car accident. performed repeatedly).
	Let \(\mathbb{P}(\text{success})=p\), \(\mathbb{P}(\text{failure})=q=1-p\).
	
	\begin{definition}
		Bernoulli r.v. \((X\sim \mathrm{Be}(p))\): \(X:=\) number of successes when the experiment is performed once.
		Then \(X\in\{0,1\}\),
		\[
		\mathbb{P}(X=1)=p,\qquad \mathbb{P}(X=0)=q=1-p.
		\]
	\end{definition}
	
	\begin{remark}
		\(\mathrm{Be}(p)\) can arise from any experiment.
		Let \(A\subseteq \Omega\). Define the indicator function:
		\[
		\mathbf{1}_A:\Omega\to \mathbb{R},\qquad
		\mathbf{1}_A(\omega)=
		\begin{cases}
			1, & \omega\in A,\\
			0, & \omega\notin A.
		\end{cases}
		\]
		Then
		\[
		\mathbb{P}(\mathbf{1}_A=0)=1-\mathbb{P}(A),
		\mathbb{P}(\mathbf{1}_A=1)=\mathbb{P}(A).
		\]
	\end{remark}
	
	\begin{definition}
		Binomial r.v. \((X\sim \mathrm{Bin}(n,p))\):
		\(X:=\) number of successes if the experiment is repeated \(n\) times.
		Then \(X\in\{0,1,2,\dots,n\}\), and for \(k=0,1,\dots,n\),
		\[
		\mathbb{P}(X=k)=\binom{n}{k}p^k(1-p)^{\,n-k}.
		\]
	\end{definition}
	
	\begin{definition}
		Geometric r.v. \((X\sim \mathrm{Geom}(p))\):
		\(X:=\) waiting time for the first success, i.e. the number of trials required until the first success.
		Then \(X\in\{1,2,\dots\}\), and for \(k=1,2,\dots\),
		\[
		\mathbb{P}(X=k)=(1-p)^{k-1}p.
		\]
	\end{definition}
	
	\begin{definition}
		Negative Binomial r.v. \((X\sim \mathrm{NB}(r,p))\):
		\(X:=\) waiting time for \(r\) successes.
		Then \(X\in\{r,r+1,\dots\}\), and for \(k=r,r+1,\dots\),
		\[
		\mathbb{P}(X=k)=\binom{k-1}{r-1}p^r(1-p)^{\,k-r}.
		\]
	\end{definition}
	
	\begin{definition}
		Poisson r.v. \((X\sim \mathrm{Poisson}(\lambda))\), \(\lambda>0\):
		for \(k=0,1,2,\dots\),
		\[
		\mathbb{P}(X=k)=\frac{\lambda^k}{k!}e^{-\lambda}.
		\]
	\end{definition}
	
	Motivation: can be regarded as an approximation of \(\mathrm{Bin}(n,p)\) in the case
	\(n\) is large, \(p\) is small, and \(np\) is moderate.
	
	\begin{definition}
		\textcolor{red}{In high school, we studied the hypergeometric distribution. Basically, given $N$ objects, and $K$ of them are regarded as success. The probability of $k$ successes when $n$ experiments are performed is given by:}
		\begin{align*}
			\mathbb{P}(X=k)
			& =
			\frac{{K\choose k}{N-K\choose n-k}}{{N\choose n}}
		\end{align*}
	\end{definition}
	
	\section{Continuous random variable}
	
	\begin{definition}
		[\textbf{Absolutely Continuous Function}]
		\textcolor{red}{A function $F(x)$ on $\mathbb{R}$ is absolutely continuous, if it admits an integrable density function $f(x)$.}
	\end{definition}
	\begin{remark}
		\textcolor{red}{If $F(x)$ is indeed absolutely continuous, then the corresponding density function $f(x)$ is not unique. Any two density functions $f_0(x),f_1(x)$ of $F(x)$ are almost everywhere the same. This analysis concept can be used to define continuous probability distribution.}
	\end{remark}
	
	For discrete r.v.,
	\[
	\mathbb{P}(a<X\le b)=F(b)-F(a).
	\]
	
	\begin{definition}
		We say a r.v. \(X\) is continuous if its c.d.f. \(F(x)=\mathbb{P}(X\le x)\) can be written as
		\[
		F(x)=\int_{-\infty}^{x} f(u)\,du
		\]
		for some integrable function \(f:\mathbb{R}\to[0,\infty)\).
		\(f\) is called probability density function (p.d.f.).
	\end{definition}
	
	Comparison with discrete:
	\[
	F(x)=\sum_{x_i\le x}\mathbb{P}(X=x_i)
	\qquad\text{vs.}\qquad
	F(x)=\int_{-\infty}^{x} f(u)\,du.
	\]
	\[
	\mathbb{P}(X=x)=F(x)-F(x-),
	\qquad
	\text{and for continuous }X,\quad \mathbb{P}(X=x)=0.
	\]
	\[
	f(x)=F'(x)\quad\text{\textcolor{red}{almost everywhere}}
	\]
	
	\begin{remark}
		\(f(x)\) is not a probability itself, but we can regard \(f(x)\,dx\) as a probability:
		\[
		\mathbb{P}(x\le X\le x+dx)=\int_{x}^{x+dx} f(u)\,du \approx f(x)\,dx.
		\]
	\end{remark}
	
	Properties: If \(X\) has p.d.f. \(f\),
	\[
	\int_{-\infty}^{\infty} f(x)\,dx=1,
	\qquad
	\mathbb{P}(X=x)=0,
	\qquad
	\mathbb{P}(a\le X\le b)=\int_{a}^{b} f(x)\,dx.
	\]
	More generally,
	\[
	\mathbb{P}(X\in B)=\int_{B} f(x)\,dx
	\quad\text{for sufficiently nice }B.
	\]
	
	Important examples of continuous r.v.s.
	
	\begin{definition}
		Uniform r.v. \((X\sim U(a,b))\).
		\(X\) is uniform on \([a,b]\) if
		\[
		F(x)=
		\begin{cases}
			0, & x<a,\\[4pt]
			\dfrac{x-a}{b-a}, & a\le x<b,\\[8pt]
			1, & x\ge b,
		\end{cases}
		\qquad
		f(x)=
		\begin{cases}
			\dfrac{1}{b-a}, & a\le x\le b,\\[8pt]
			0, & \text{otherwise}.
		\end{cases}
		\]
	\end{definition}
	
	An application: Inverse sampling transform.
	Given a distribution \(F\), how to generate a r.v. \(X\) with c.d.f. \(F\)?
	Suppose we can generate \(U\sim U(0,1)\).
	If \(F\) is continuous and invertible (\textcolor{red}{a homeomorphism}), set
	\[
	X=F^{-1}(U).
	\]
	Claim: \(X\) has c.d.f. \(F\). Indeed,
	\[
	\mathbb{P}(X\le x)
	=\mathbb{P}(F^{-1}(U)\le x)
	=\mathbb{P}(U\le F(x))
	=F(x).
	\]
	
	\begin{definition}
		Exponential distribution \((X\sim \mathrm{Exp}(\lambda))\), \(\lambda>0\).
		\(X\) is exponential with parameter \(\lambda\) if
		\[
		F(x)=
		\begin{cases}
			0, & x<0,\\
			1-e^{-\lambda x}, & x\ge 0,
		\end{cases}
		\qquad
		f(x)=
		\begin{cases}
			\lambda e^{-\lambda x}, & x\ge 0,\\
			0, & x<0.
		\end{cases}
		\]
	\end{definition}
	
	Motivation:
	Poisson process vs Binomial process.
	\[
	\mathrm{Exp}(\lambda):\ \text{waiting time for the first success in a Poisson process}, \]
	\[
	\mathrm{Geom}(p):\ \text{waiting time for the first success in a Binomial process}.
	\]
	
	\begin{definition}
		Gamma r.v. \((X\sim \mathrm{Gamma}(\alpha,\lambda))\), \(\alpha,\lambda>0\):
		\[
		f(x)=
		\begin{cases}
			\dfrac{\lambda^{\alpha}}{\Gamma(\alpha)}x^{\alpha-1}e^{-\lambda x}, & x>0,\\[10pt]
			0, & x\le 0,
		\end{cases}
		\qquad
		\Gamma(\alpha)=\int_{0}^{\infty} e^{-t}t^{\alpha-1}\,dt.
		\]
	\end{definition}
	
	Motivation: if \(\alpha\) is an integer, \(X\) is waiting time for the \(\alpha\)-th success during a Poisson process.
	
	\begin{definition}
		Gaussian (Normal) r.v. \((X\sim N(\mu,\sigma^2))\), \(\sigma>0\):
		\[
		f(x)=\frac{{1}}{\sqrt{2\pi}\sigma}\exp\!\left(-\frac{(x-\mu)^2}{2\sigma^2}\right),
		\qquad x\in\mathbb{R}.
		\]
		Standard normal: \(\mu=0\), \(\sigma^2=1\).
		\[
		\varphi(x)=\frac{1}{\sqrt{2\pi}}e^{-x^2/2},
		\qquad
		\Phi(x)=\int_{-\infty}^{x}\varphi(u)\,du.
		\]
		Importance: central limit theorem.
	\end{definition}
	
	\begin{proposition}
		(Linear transform of Gaussian is still Gaussian.)
		If \(X\sim N(\mu,\sigma^2)\) and \(Y=aX+b\) with \(a\ne 0\), then
		\[
		Y\sim N(a\mu+b,\ a^2\sigma^2).
		\]
	\end{proposition}
	
	\begin{remark}
		If \(X\sim N(0,1)\) and \(Y=\sigma X+\mu\), then
		\[
		Y\sim N(\mu,\sigma^2).
		\]
	\end{remark}
	
	
	\section{Expectation of R.V. /distribution}
	
	
	
	Motivation: theoretical approximation of numerical/empirical average.
	
	Let \(X_1, X_2, \dots, X_N\) be numerical outcomes of \(N\) repetitions of the same experiment.
	
	Empirical average:
	\[
	m=\frac{1}{N}\sum_{i=1}^{N}X_i.
	\]
	
	When \(N\) is large, intuition suggests that the number of occurrences of a specific value
	\(x_k \in \mathrm{Range}(X)\) among \(X_1,\dots,X_N\) should be about
	\(N\mathbb{P}(X=x_k)\). Hence
	\[
	m=\frac{1}{N}\sum_{i=1}^{N}X_i
	\approx \frac{1}{N}\sum_{k}\bigl(N\mathbb{P}(X=x_k)\bigr)x_k
	=\sum_k x_k\,\mathbb{P}(X=x_k).
	\]
	
	\begin{definition}
		(discrete case) Let \(X\) be discrete with p.m.f.\ \(p_X(x)=\mathbb{P}(X=x)\).
		Its expectation is
		\[
		\mathbb{E}X := \sum_x x\,p_X(x),
		\]
		(the mean / weighted average).
	\end{definition}
	
	\begin{definition}
		(continuous case) Let \(X\) be continuous with p.d.f.\ \(f_X(x)\).
		Its expectation is
		\[
		\mathbb{E}X := \int_{-\infty}^{\infty} x f_X(x)\,dx.
		\]
	\end{definition}
	
	\begin{remark}
		\textcolor{red}{Review of some introductory analysis: Let $a_n$ be a sequence of nonnegative real numbers, the value of $\sum_na_n$ doesn't depend on the order of summation. Let $a_n$ be a sequence of not necessarily nonnegative real numbers. If $\sum_n|a_n|$ converges, then both the nonpositive part and the nonnegative part of $\sum_na_n$ converge by comparison test. In this case, it follows from the $\lim$ is linear that $\sum_na_n$ converge. We say that $\sum_na_n$ converges absolutely. If $\sum_n|a_n|$ diverges, but $\sum_na_n$ converges, then $\sum_na_n$ converges conditionally. Riemann arrangement theorem tells that for all real number $A$, for some bijection $\varphi$ on $\mathbb{N}$, $\sum_na_{\varphi(n)}$ converges to $A$.}
	\end{remark}
	
	\begin{remark}
		\textcolor{red}{The probability density function $f_X(x)$ is nonnegative, so the cumulative distribution function $F_X(x)=\int_{-\infty}^xf_X(y)\mathrm{d}y$, or $F_X(x)=\sum_{y=-\infty}^{y=x}f_X(y)$ doesn't depend on the reordering of the subscript. However, when computing the expectation $\mathbb{E}X$ of a random variable, the factor $x$ is introduced, and we notice that $x$ is not necessarily nonnegative, so the order of summation matters here. Professor Bao mentioned that we only define expectation when the integral/sum converges absolutely.}
	\end{remark}
	
	\begin{example}
		\(X\sim \mathrm{Be}(p)\). Then
		\[
		\mathbb{E}X=\sum_{x\in\{0,1\}} x\,\mathbb{P}(X=x)
		=0\cdot\mathbb{P}(X=0)+1\cdot\mathbb{P}(X=1)
		=0\cdot(1-p)+1\cdot p
		=p.
		\]
	\end{example}
	
	\begin{example}
		\(X\sim \mathrm{Bin}(n,p)\). Then
		\[
		\mathbb{E}X=\sum_{k=0}^{n} k\,\mathbb{P}(X=k)
		=\sum_{k=0}^{n} k \binom{n}{k}p^k(1-p)^{n-k}.
		\]
		Use
		\[
		k\binom{n}{k}=n\binom{n-1}{k-1}\qquad (k=1,2,\dots,n),
		\]
		so
		\[
		\mathbb{E}X
		=\sum_{k=1}^{n} n\binom{n-1}{k-1}p^k(1-p)^{n-k}
		=np\sum_{k=1}^{n} \binom{n-1}{k-1}p^{k-1}(1-p)^{(n-1)-(k-1)}.
		\]
		Let \(j=k-1\). Then \(j=0,1,\dots,n-1\), and
		\[
		\mathbb{E}X
		=np\sum_{j=0}^{n-1} \binom{n-1}{j}p^{j}(1-p)^{(n-1)-j}
		=np\,(p+(1-p))^{n-1}
		=np.
		\]
	\end{example}
	
	\begin{example}
		\(X\sim N(\mu,\sigma^2)\). Then
		\[
		f_X(x)=\frac{1}{\sqrt{2\pi}\sigma}\exp\!\left(-\frac{(x-\mu)^2}{2\sigma^2}\right),
		\qquad x\in\mathbb{R}.
		\]
		Compute
		\[
		\mathbb{E}X=\int_{-\infty}^{\infty} x f_X(x)\,dx.
		\]
		Let \(y=\frac{x-\mu}{\sigma}\), i.e. \(x=\sigma y+\mu\), \(dx=\sigma\,dy\). Then
		\[
		\mathbb{E}X
		=\int_{-\infty}^{\infty} (\sigma y+\mu)\frac{1}{\sqrt{2\pi}\sigma}e^{-y^2/2}\,\sigma\,dy
		=\int_{-\infty}^{\infty} (\sigma y+\mu)\frac{1}{\sqrt{2\pi}}e^{-y^2/2}\,dy.
		\]
		Split the integral:
		\[
		\mathbb{E}X
		=\sigma\int_{-\infty}^{\infty} y\frac{1}{\sqrt{2\pi}}e^{-y^2/2}\,dy
		+\mu\int_{-\infty}^{\infty} \frac{1}{\sqrt{2\pi}}e^{-y^2/2}\,dy.
		\]
		The first integral is \(0\) since the integrand is odd:
		\[
		\int_{-\infty}^{\infty} y\frac{1}{\sqrt{2\pi}}e^{-y^2/2}\,dy=0.
		\]
		The second integral is \(1\) since it is the total probability of \(N(0,1)\):
		\[
		\int_{-\infty}^{\infty} \frac{1}{\sqrt{2\pi}}e^{-y^2/2}\,dy=1.
		\]
		Therefore
		\[
		\mathbb{E}X=\mu.
		\]
	\end{example}
	
	Expectation of a function of a r.v.
	
	Discrete case: \(X\) a discrete r.v. with p.m.f. \(p_X(x)\).
	Let \(g:\mathbb{R}\to\mathbb{R}\). Define \(Y=g(X)\).
	By definition of \(\mathbb{E}(\cdot)\), for \(Y=g(X)\), we can identify the range of \(Y\)
	and its p.m.f. \(p_Y(\cdot)\), then
	\[
	\mathbb{E}g(X)=\mathbb{E}Y=\sum_b b\,p_Y(b).
	\]
	
	\begin{proposition}
		If \(X\) is discrete, \textcolor{red}{and $g:\mathbb{R}\rightarrow\mathbb{R}$ is a ``nice" function (again should clarify what it means, I think Prof. Bao mentioned that one condition for this proposition is that $g$ is continuous and integrable. I may go wrong, please check the textbook)}, then
		\[
		\mathbb{E}g(X)=\sum_x g(x)\,p_X(x).
		\]
	\end{proposition}
	
	Continuous: \(X\) is continuous with p.d.f. \(f_X(x)\), \(g:\mathbb{R}\to\mathbb{R}\).
	\[
	\mathbb{E}g(X)=\int_{-\infty}^{\infty} g(x) f_X(x)\,dx.
	\]
	
	e.g. \(k\)-th moment: take \(g(x)=x^k\). Then
	\[
	\mathbb{E}(X^k)=
	\begin{cases}
		\displaystyle\sum_x x^k\,p_X(x), & \text{discrete},\\[10pt]
		\displaystyle\int_{-\infty}^{\infty} x^k f_X(x)\,dx, & \text{continuous}.
	\end{cases}
	\]
	
	\(k\)-th central moment: take \(g(x)=(x-\mathbb{E}X)^k\). Then
	\[
	\mathbb{E}\big[(X-\mathbb{E}X)^k\big]=
	\begin{cases}
		\displaystyle\sum_x (x-\mathbb{E}X)^k\,p_X(x), & \text{discrete},\\[10pt]
		\displaystyle\int_{-\infty}^{\infty} (x-\mathbb{E}X)^k f_X(x)\,dx, & \text{continuous}.
	\end{cases}
	\]
	
	\(k=2\). Variance:
	\[
	\mathrm{Var}(X):=\mathbb{E}\big[(X-\mathbb{E}X)^2\big],
	\]
	measure of dispersion/randomness of \(X\).
	Standard deviation:
	\[
	\mathrm{SD}(X)=\sqrt{\mathrm{Var}(X)}.
	\]
	\begin{remark}
	    \textcolor{red}{It is common to denote the expectation of $X$:}
        \begin{align*}
            \textcolor{red}{\mathbb{E}X}
        \end{align*}
        \textcolor{red}{As:}
        \begin{align*}
            \textcolor{red}{\mathbb{E}[X]}
        \end{align*}
        \textcolor{red}
        {
            The square bracket notation here means that $X$ is a function,
            and $\mathbb{E}$ is a linear functional.
        }
	\end{remark}
	\begin{proposition}
		\[
		\mathrm{Var}(X)=\mathbb{E}\big[(X-\mathbb{E}X)^2\big]=\mathbb{E}(X^2)-(\mathbb{E}X)^2.
		\]
	\end{proposition}
	
	\begin{proposition}
		For \(a,b\in\mathbb{R}\),
		\[
		\mathrm{Var}(aX+b)=a^2\mathrm{Var}(X).
		\]
	\end{proposition}
	
	
	\section{Random vector/ jointly distributed random variables}
	
	Random vector (jointly distributed r.v.'s) is a vector-valued random variable.
	It is used to describe multiple numerical quantities associated to the same outcome \(\omega\).
	
	\begin{example}
		(Basketball players.)
		Suppose there are \(5\) candidates, and we pick one uniformly at random.
		Let the sample space be \(\Omega=\{1,2,3,4,5\}\) with \(\mathbb{P}(\{\omega\})=\frac{1}{5}\).
		Define height \(H(\omega)\) and weight \(W(\omega)\) by
		\[
		\begin{array}{c|cc}
			\omega & H(\omega) & W(\omega)\\\hline
			1 & 180 & 69\\
			2 & 182 & 75\\
			3 & 175 & 65\\
			4 & 174 & 70\\
			5 & 170 & 68
		\end{array}
		\]
		Then \((H,W)\) is a random vector.
		
		Probability that a randomly picked one is qualified, i.e. \(H\ge 175\) and \(W\ge 70\):
		\[
		\mathbb{P}(H\ge 175,\ W\ge 70)
		=\mathbb{P}\big(\{\omega:\ H(\omega)\ge 175\}\cap\{\omega:\ W(\omega)\ge 70\}\big).
		\]
		Compute
		\[
		\{\omega:\ H(\omega)\ge 175\}=\{1,2,3\},
		\qquad
		\{\omega:\ W(\omega)\ge 70\}=\{2,4\},
		\]
		so the intersection is \(\{2\}\). Hence
		\[
		\mathbb{P}(H\ge 175,\ W\ge 70)=\mathbb{P}(\{2\})=\frac{1}{5}.
		\]
		Also,
		\[
		\mathbb{P}(H\ge 175)=\frac{3}{5},\qquad
		\mathbb{P}(W\ge 70)=\frac{2}{5}.
		\]
		In this example,
		\[
		\mathbb{P}(H\ge 175,\ W\ge 70)=\frac{1}{5}\ne \frac{3}{5}\cdot\frac{2}{5},
		\]
		so \(H\) and \(W\) are not independent.
	\end{example}
	
	More generally, for any set \(B\subseteq \mathbb{R}^2\), we may ask \(\mathbb{P}((X,Y)\in B)\).
	This probability is determined by the joint distribution of \((X,Y)\).
	
	\begin{definition}
		Joint cdf of \((X,Y)\) is
		\[
		F_{X,Y}(x,y)=\mathbb{P}(X\le x,\ Y\le y)
		=\mathbb{P}\big((X,Y)\in (-\infty,x]\times(-\infty,y]\big).
		\]
	\end{definition}
	
	More generally, for \(X_1,\dots,X_n\),
	\[
	F_{X_1,\dots,X_n}(x_1,\dots,x_n)
	=\mathbb{P}(X_1\le x_1,\dots, X_n\le x_n).
	\]
	
	\begin{remark}
		From \(F_X(x)\) and \(F_Y(y)\), we will not be able to get \(F_{X,Y}(x,y)\) in general.
	\end{remark}
	
	Marginal cdf:
	\[
	F_X(x)=\lim_{y\to +\infty}F_{X,Y}(x,y),
	\qquad
	F_Y(y)=\lim_{x\to +\infty}F_{X,Y}(x,y).
	\]
	
	We say \((X,Y)\) is jointly discrete if it takes values on a finite or countably infinite set.
	
	\begin{definition}
		Joint p.m.f.\ of \((X,Y)\) is
		\[
		p_{X,Y}(x,y)=\mathbb{P}(X=x,\ Y=y).
		\]
	\end{definition}
	
	Marginal p.m.f.'s:
	\[
	p_X(x)=\sum_{y} p_{X,Y}(x,y),
	\qquad
	p_Y(y)=\sum_{x} p_{X,Y}(x,y).
	\]
	
	Joint cdf in terms of joint p.m.f.:
	\[
	F_{X,Y}(x,y)=\sum_{u\le x}\sum_{v\le y} p_{X,Y}(u,v).
	\]
	
	(Recover joint p.m.f.\ from joint cdf; lattice-valued case.)
	If \(X,Y\) take integer values, then
	\[
	p_{X,Y}(x,y)
	=F_{X,Y}(x,y)-F_{X,Y}(x-1,y)-F_{X,Y}(x,y-1)+F_{X,Y}(x-1,y-1).
	\]
	
	Probability of a rectangle.
	For \(x_1<x_2\) and \(y_1<y_2\),
	\[
	\mathbb{P}(x_1<X\le x_2,\ y_1<Y\le y_2)
	=F_{X,Y}(x_2,y_2)-F_{X,Y}(x_1,y_2)-F_{X,Y}(x_2,y_1)+F_{X,Y}(x_1,y_1).
	\]
	
	We say \(X\) and \(Y\) are jointly continuous if there exists \(f_{X,Y}:\mathbb{R}^2\to[0,\infty)\)
	such that
	\[
	F_{X,Y}(x,y)=\mathbb{P}(X\le x,\ Y\le y)
	=\int_{-\infty}^{x}\int_{-\infty}^{y} f_{X,Y}(u,v)\,dv\,du.
	\]
	In this case, \(f_{X,Y}\) is called the joint p.d.f.
	
	\begin{remark}
		If both \(X\) and \(Y\) are continuous, they may not be jointly continuous.
		e.g.\ \(X\sim U(0,1)\), let \(Y=X\). Then \((X,Y)\) lies on the line \(y=x\), so no joint p.d.f.\ exists.
	\end{remark}
	
	If \(F_{X,Y}\) is smooth enough, then
	\[
	f_{X,Y}(x,y)=\frac{\partial^2}{\partial x\,\partial y}F_{X,Y}(x,y).
	\]
	
	Marginal p.d.f.:
	\[
	f_X(x)=\int_{-\infty}^{\infty} f_{X,Y}(x,y)\,dy,
	\qquad
	f_Y(y)=\int_{-\infty}^{\infty} f_{X,Y}(x,y)\,dx.
	\]
	
	\begin{proposition}
		Given \(X,Y\) and \(g:\mathbb{R}^2\to\mathbb{R}\),
		\[
		\mathbb{E}[g(X,Y)]
		=
		\begin{cases}
			\displaystyle\sum_x\sum_y g(x,y)\,p_{X,Y}(x,y), & \text{(jointly discrete)},\\[12pt]
			\displaystyle\int_{-\infty}^{\infty}\int_{-\infty}^{\infty}
			g(x,y)\,f_{X,Y}(x,y)\,dx\,dy, & \text{(jointly continuous)}.
		\end{cases}
		\]
	\end{proposition}
	
	\begin{example}
		Let \(g(x,y)=x+y\). In the jointly continuous case,
		\[
		\mathbb{E}[X+Y]
		=\int_{-\infty}^{\infty}\int_{-\infty}^{\infty} (x+y)f_{X,Y}(x,y)\,dx\,dy
		\]
		\[
		=\int\!\!\int x f_{X,Y}(x,y)\,dx\,dy+\int\!\!\int y f_{X,Y}(x,y)\,dx\,dy
		=\mathbb{E}X+\mathbb{E}Y.
		\]
	\end{example}
	
	\begin{proposition}
		(Linearity of expectation.)
		For any \(a_1,\dots,a_n\in\mathbb{R}\),
		\[
		\mathbb{E}\!\left[\sum_{i=1}^{n} a_i X_i\right]
		=\sum_{i=1}^{n} a_i \mathbb{E}[X_i].
		\]
	\end{proposition}
	
	\begin{example}
		If \(X\sim \mathrm{Bin}(n,p)\), then \(X=\sum_{i=1}^{n}Y_i\) where \(Y_i\sim \mathrm{Be}(p)\) i.i.d.
		Therefore
		\[
		\mathbb{E}X=\mathbb{E}\!\left[\sum_{i=1}^{n}Y_i\right]
		=\sum_{i=1}^{n}\mathbb{E}[Y_i]
		=\sum_{i=1}^{n}p
		=np.
		\]
	\end{example}
	
	
	
	\begin{example}
		
		Coupon collection problem: \(N\) different types of card and each time one obtains a card,
		which is equally likely to be any of \(N\) types.
		
		Q1: Expected number of different types contained in \(n\) cards.
		
		Let \(X\) be the number of different types contained in \(n\) cards.
		Write
		\[
		X=X_1+X_2+\cdots+X_N,
		\]
		where
		\[
		X_i=
		\begin{cases}
			1, & \text{if at least one type \(i\) is among \(n\) cards},\\
			0, & \text{otherwise}.
		\end{cases}
		\]
		Then
		\[
		\mathbb{E}X=\sum_{i=1}^{N}\mathbb{E}X_i=\sum_{i=1}^{N}\mathbb{P}(X_i=1).
		\]
		Also
		\[
		\mathbb{P}(X_i=1)=1-\mathbb{P}(X_i=0)
		=1-\left(\frac{N-1}{N}\right)^{n}.
		\]
		Hence
		\[
		\mathbb{E}X
		=\sum_{i=1}^{N}\left(1-\left(\frac{N-1}{N}\right)^n\right)
		=N\left(1-\left(\frac{N-1}{N}\right)^n\right).
		\]
		
		Q2: Expected number of cards one need to collect until one has a complete set.
		
		Let \(Y\) be the number of cards needed for a full collection.
		Decompose
		\[
		Y=Y_0+Y_1+\cdots+Y_{N-1},
		\]
		where \(Y_i\) is the number of cards needed to be collected after \(i\) distinct types have already been obtained,
		in order to get another new type.
		
		When \(i\) distinct types are obtained, the probability the next card is a new type is
		\[
		p_i=\frac{N-i}{N}.
		\]
		Hence
		\[
		Y_i\sim \mathrm{Geom}(p_i),
		\qquad
		\mathbb{E}Y_i=\frac{1}{p_i}=\frac{N}{N-i}.
		\]
		Therefore
		\[
		\mathbb{E}Y=\sum_{i=0}^{N-1}\mathbb{E}Y_i
		=\sum_{i=0}^{N-1}\frac{N}{N-i}
		=N\sum_{k=1}^{N}\frac{1}{k}.
		\]
		
	\end{example}
	
	
	
	\section{Independence of random variables}
	We say \(X\) and \(Y\) are independent denoted \(X\perp\!\!\!\perp Y\) (\textcolor{red}{This notation is different from uncorrelated, be careful!}) if for any Borel sets \(A,B\subseteq\mathbb{R}\),
	\[
	\mathbb{P}(X\in A,\ Y\in B)=\mathbb{P}(X\in A)\,\mathbb{P}(Y\in B).
	\]
	Otherwise, we say \(X\) and \(Y\) are dependent.
	
	
	\begin{definition}
		(Independence via cdf.)
		We say \(X\) and \(Y\) are independent if for any \(x,y\in\mathbb{R}\),
		\[
		\mathbb{P}(X\le x,\ Y\le y)=\mathbb{P}(X\le x)\,\mathbb{P}(Y\le y),
		\qquad\text{i.e.}\qquad
		F_{X,Y}(x,y)=F_X(x)F_Y(y).
		\]
	\end{definition}
	
	\begin{remark}
		(Typical scenarios.) Coin tossing, die rolling, and more generally independent repeated trials
		often produce independent random variables.
	\end{remark}
	
	\begin{remark}
		(Generalization.)
		Random variables \(X_1,\dots,X_n\) are independent if for any \(x_1,\dots,x_n\in\mathbb{R}\),
		\[
		F_{X_1,\dots,X_n}(x_1,\dots,x_n)=\prod_{i=1}^{n} F_{X_i}(x_i).
		\]
		Equivalently, for any Borel sets \(A_1,\dots,A_n\subseteq\mathbb{R}\),
		\[
		\mathbb{P}(X_1\in A_1,\dots, X_n\in A_n)=\prod_{i=1}^{n}\mathbb{P}(X_i\in A_i).
		\]
	\end{remark}
	
	\textbf{Alternative characterizations.}
	
	(i) If \((X,Y)\) is jointly discrete, then \(X\perp\!\!\!\perp Y\) iff for all \(x,y\),
	\[
	p_{X,Y}(x,y)=p_X(x)\,p_Y(y).
	\]
	
	(ii) If \((X,Y)\) is jointly continuous, then \(X\perp\!\!\!\perp Y\) iff for all \(x,y\),
	\[
	f_{X,Y}(x,y)=f_X(x)\,f_Y(y).
	\]
	
	\begin{proposition}
		If \(X\perp\!\!\!\perp Y\), then for any functions \(g,h:\mathbb{R}\to\mathbb{R}\) such that expectations exist,
		\[
		\mathbb{E}\big[g(X)h(Y)\big]=\mathbb{E}[g(X)]\,\mathbb{E}[h(Y)].
		\]
	\end{proposition}
	
	\begin{remark}
		(Warning: expectation may not exist.)
		For heavy-tailed distributions (e.g.\ Cauchy), certain moments may be infinite.
		For example, if \(X\) is standard Cauchy with density
		\[
		f_X(x)=\frac{1}{\pi(1+x^2)},\qquad x\in\mathbb{R},
		\]
		then
		\[
		\mathbb{E}[X^2]=\int_{-\infty}^{\infty} x^2 \frac{1}{\pi(1+x^2)}\,dx=\infty,
		\]
		so \(\mathrm{Var}(X)\) is not defined.
	\end{remark}
	
	\textbf{Variance of sum and covariance.}
	Let \(X_1,\dots,X_n\) be random variables with finite second moments.
	Then
	\[
	\mathrm{Var}\!\left(\sum_{i=1}^{n} X_i\right)
	=\mathbb{E}\!\left[\left(\sum_{i=1}^{n} (X_i-\mathbb{E}X_i)\right)^2\right]
	=\sum_{i=1}^{n}\mathrm{Var}(X_i)+2\sum_{1\le i<j\le n}\mathrm{Cov}(X_i,X_j).
	\]
	
	\begin{definition}
		(Covariance.)
		For random variables \(X,Y\) with finite second moments,
		\[
		\mathrm{Cov}(X,Y)=\mathbb{E}\big[(X-\mathbb{E}X)(Y-\mathbb{E}Y)\big]
		=\mathbb{E}[XY]-\mathbb{E}X\,\mathbb{E}Y.
		\]
	\end{definition}
	
	Note that
	\[
	\mathrm{Var}(X)=\mathrm{Cov}(X,X).
	\]
	
	\begin{remark}
		If \(X\perp\!\!\!\perp Y\), then
		\[
		\mathrm{Cov}(X,Y)=0.
		\]
		Therefore, if \(X_1,\dots,X_n\) are independent and have finite second moments, then
		\[
		\mathrm{Var}\!\left(\sum_{i=1}^{n} X_i\right)=\sum_{i=1}^{n}\mathrm{Var}(X_i).
		\]
	\end{remark}
	
	\textbf{Correlated vs.\ uncorrelated.}
	We say \(X\) and \(Y\) are uncorrelated if
	\[
	\mathrm{Cov}(X,Y)=0
	\qquad\text{(equivalently, } \mathbb{E}[XY]=\mathbb{E}X\,\mathbb{E}Y\text{)}.
	\]
	
	\begin{remark}
		Uncorrelated $X\perp Y$ does \emph{not} imply independent $X\perp\!\!\!\perp Y$ in general.
	\end{remark}
	
	\begin{example}
		(Uncorrelated but dependent.)
		Let \(X\) be a discrete r.v.\ such that
		\[
		\mathbb{P}(X=-1)=\mathbb{P}(X=0)=\mathbb{P}(X=1)=\frac{1}{3}.
		\]
		Define
		\[
		Y=
		\begin{cases}
			0, & X\ne 0,\\
			1, & X=0.
		\end{cases}
		\]
		Then \(Y\) is a function of \(X\), so \(X\) and \(Y\) are dependent.
		But
		\[
		\mathbb{E}X=0,
		\qquad
		\mathbb{E}Y=\frac{1}{3},
		\qquad
		\mathbb{E}[XY]=0,
		\]
		so
		\[
		\mathrm{Cov}(X,Y)=\mathbb{E}[XY]-\mathbb{E}X\,\mathbb{E}Y=0,
		\]
		i.e.\ they are uncorrelated.
		Also,
		\[
		\mathbb{P}(X=0,\ Y=1)=\frac{1}{3}\ne \frac{1}{3}\cdot\frac{1}{3}=\mathbb{P}(X=0)\mathbb{P}(Y=1),
		\]
		so they are not independent.
	\end{example}
	
	
	
	\section{Distribution of sum of r.v.}
	
	Let \(X,Y\) be two random variables and define \(Z=X+Y\).
	We want to determine the distribution (p.m.f./p.d.f./c.d.f.) of \(Z\) from that of \(X\) and \(Y\).
	
	\subsection{Discrete case}
	
	Assume \((X,Y)\) is jointly discrete with joint p.m.f.
	\[
	p_{X,Y}(x,y)=\mathbb{P}(X=x,\ Y=y).
	\]
	Then for any \(z\in\mathbb{R}\) (in practice, \(z\) ranges over the lattice support),
	\begin{align*}
		p_Z(z)
		&=\mathbb{P}(Z=z)
		=\mathbb{P}(X+Y=z) \\
		&=\sum_{x}\mathbb{P}(X=x,\ Y=z-x)
		=\sum_x p_{X,Y}(x,z-x).
	\end{align*}
	
	\begin{remark}
		If \(X\perp\!\!\!\perp Y\), then \(p_{X,Y}(x,y)=p_X(x)p_Y(y)\), so
		\[
		p_Z(z)=\sum_x p_X(x)\,p_Y(z-x).
		\]
		This is the \textbf{convolution} of p.m.f.'s:
		\[
		p_Z=p_X * p_Y.
		\]
	\end{remark}
	
	\begin{example}
		Let \(X\sim \mathrm{Be}(p)\), \(Y\sim \mathrm{Be}(q)\), independent.
		Then \(Z=X+Y\in\{0,1,2\}\), and
		\[
		\mathbb{P}(Z=0)=(1-p)(1-q),\quad
		\mathbb{P}(Z=2)=pq,
		\]
		\[
		\mathbb{P}(Z=1)=p(1-q)+(1-p)q=p+q-2pq.
		\]
	\end{example}
	
	\subsection{Continuous case}
	
	Assume \((X,Y)\) is jointly continuous with joint p.d.f. \(f_{X,Y}(x,y)\).
	Let \(Z=X+Y\). First write the c.d.f.:
	\begin{align*}
		F_Z(z)
		&=\mathbb{P}(Z\le z)
		=\mathbb{P}(X+Y\le z)\\
		&=\iint_{\{(x,y):\,x+y\le z\}} f_{X,Y}(x,y)\,dx\,dy.
	\end{align*}
	If we integrate in the order \(dx\) then \(dy\), we get
	\[
	F_Z(z)=\int_{-\infty}^{\infty}\left(\int_{-\infty}^{z-y} f_{X,Y}(x,y)\,dx\right)dy.
	\]
	
	If \(F_Z\) is differentiable in \(z\) (under mild regularity assumptions), the p.d.f. of \(Z\) is
	\begin{align*}
		f_Z(z)
		&=\frac{d}{dz}F_Z(z)
		=\int_{-\infty}^{\infty} f_{X,Y}(z-y,y)\,dy.
	\end{align*}
	
	\begin{remark}
		If \(X\perp\!\!\!\perp Y\), then \(f_{X,Y}(x,y)=f_X(x)f_Y(y)\), hence
		\[
		f_Z(z)=\int_{-\infty}^{\infty} f_X(z-y)\,f_Y(y)\,dy.
		\]
		This is the \textbf{convolution} of p.d.f.'s:
		\[
		f_Z=f_X * f_Y.
		\]
	\end{remark}
	
	\begin{example}
		If \(X,Y\) are i.i.d.\ \(U(0,1)\), then
		\[
		f_X(x)=f_Y(x)=
		\begin{cases}
			1, & 0\le x\le 1,\\
			0, & \text{otherwise}.
		\end{cases}
		\]
		For \(z\in\mathbb{R}\),
		\begin{align*}
			f_Z(z)
			&=\int_{-\infty}^{\infty} \mathbf{1}_{[0,1]}(z-y)\mathbf{1}_{[0,1]}(y)\,dy\\
			&=\int_{-\infty}^{\infty} \mathbf{1}_{[0,1]}(y)\mathbf{1}_{[0,1]}(z-y)\,dy
			=\int_{\max\{0,z-1\}}^{\min\{1,z\}} 1\,dy.
		\end{align*}
		Therefore
		\[
		f_Z(z)=
		\begin{cases}
			0, & z<0,\\
			z, & 0\le z\le 1,\\
			2-z, & 1< z\le 2,\\
			0, & z>2.
		\end{cases}
		\]
		So \(Z=X+Y\) has the triangular distribution on \([0,2]\).
	\end{example}
	

	
	\begin{proposition}
		If \(X\sim N(\mu_1,\sigma_1^2)\) and \(Y\sim N(\mu_2,\sigma_2^2)\) are independent, then
		\[
		X+Y \sim N(\mu_1+\mu_2,\ \sigma_1^2+\sigma_2^2).
		\]
		More generally, if \(X_1,\dots,X_n\) are independent with \(X_i\sim N(\mu_i,\sigma_i^2)\), then
		\[
		\sum_{i=1}^n X_i \sim N\!\left(\sum_{i=1}^n \mu_i,\ \sum_{i=1}^n \sigma_i^2\right).
		\]
	\end{proposition}
	
	\begin{remark}
		A special case: if \(X_1,\dots,X_n\) are i.i.d.\ \(N(\mu,\sigma^2)\), then
		\[
		\overline{X}_n=\frac{1}{n}\sum_{i=1}^n X_i \sim N\!\left(\mu,\ \frac{\sigma^2}{n}\right).
		\]
	\end{remark}
	
	\begin{note}
	Convolution is often computationally complicated. Later we will introduce tools (e.g.\ moment generating functions / characteristic functions) to track distributions of sums more easily, especially under independence.
	\end{note}
	
	
	\section{Moment generating function (Laplace transform)}
	
	\textbf{Motivation.}
Convolution is often complicated. For sums of independent r.v.'s, m.g.f. turns convolution into multiplication.
	
	\begin{definition}
		(Moment generating function.)
		Let \(X\) be a random variable. Its moment generating function (m.g.f.) is
		\[
		\phi_X(t):=\mathbb{E}\big[e^{tX}\big],
		\]
		for all \(t\in\mathbb{R}\) such that the expectation exists (possibly only for \(t\) in an open interval containing \(0\)).
	\end{definition}
	
	\begin{remark}
		(Discrete vs.\ continuous forms.)
		\begin{itemize}[leftmargin=2em]
			\item If \(X\) is discrete with p.m.f.\ \(p_X(x)\), then
			\[
			\phi_X(t)=\sum_x e^{tx}p_X(x).
			\]
			\item If \(X\) is continuous with p.d.f.\ \(f_X(x)\), then
			\[
			\phi_X(t)=\int_{-\infty}^{\infty} e^{tx} f_X(x)\,dx.
			\]
		\end{itemize}
	\end{remark}
	
	\begin{example}
		If \(X\sim \mathrm{Exp}(\lambda)\) with density \(f_X(x)=\lambda e^{-\lambda x}\mathbf{1}_{\{x\ge 0\}}\), then
		\[
		\phi_X(t)=\int_{0}^{\infty} e^{tx}\lambda e^{-\lambda x}\,dx
		=\lambda\int_{0}^{\infty} e^{-(\lambda-t)x}\,dx
		=
		\begin{cases}
			\dfrac{\lambda}{\lambda-t}, & t<\lambda,\\[8pt]
			\text{does not exist}, & t\ge \lambda.
		\end{cases}
		\]
	\end{example}
	
	\begin{example}
		\textcolor{red}{Cauchy warning.}
		If \(X\) is standard Cauchy with density \(f_X(x)=\dfrac{1}{\pi(1+x^2)}\), then
		\[
		\phi_X(t)=\mathbb{E}[e^{tX}]
		\]
		does \emph{not} exist for any \(t\ne 0\) (the tails are too heavy).
		So m.g.f.\ is not always a valid tool.
	\end{example}
	

	
	\begin{proposition}
		(Independence \(\Rightarrow\) product rule.)
		If \(X\perp\!\!\!\perp Y\) and \(\phi_X(t),\phi_Y(t)\) are finite for some \(t\), then
		\[
		\phi_{X+Y}(t)=\phi_X(t)\phi_Y(t).
		\]
	\end{proposition}
	
	\begin{proof}
		\[
		\phi_{X+Y}(t)=\mathbb{E}\big[e^{t(X+Y)}\big]
		=\mathbb{E}\big[e^{tX}e^{tY}\big]
		=\mathbb{E}[e^{tX}]\,\mathbb{E}[e^{tY}]
		=\phi_X(t)\phi_Y(t),
		\]
		using the factorization property \(\mathbb{E}[g(X)h(Y)]=\mathbb{E}g(X)\,\mathbb{E}h(Y)\) for independent r.v.'s.
	\end{proof}
	
	\begin{proposition}
		(Distribution determined by m.g.f.)
		If \(\phi_X(t)\) and \(\phi_Y(t)\) are finite on some open interval containing \(0\) and
		\[
		\phi_X(t)=\phi_Y(t)\quad \text{for all \(t\) in that interval},
		\]
		then \(X\) and \(Y\) have the same distribution.
	\end{proposition}
	
	\begin{remark}
	This is a uniqueness theorem (can be proved using analytic continuation / inverse Laplace transform ideas).
		It requires existence on an interval around \(0\); without that, m.g.f.\ may fail to exist and cannot characterize the distribution.
	\end{remark}
	
	\begin{proposition}
		(Moments via m.g.f.)
		Suppose \(\phi_X(t)\) is finite in a neighborhood of \(0\). Then for each \(n\in\mathbb{N}\),
		\[
		\phi_X^{(n)}(0)=\mathbb{E}[X^n].
		\]
		In particular,
		\[
		\phi_X(0)=1,\qquad \phi_X'(0)=\mathbb{E}X,\qquad \phi_X''(0)=\mathbb{E}[X^2],
		\qquad \mathrm{Var}(X)=\phi_X''(0)-\big(\phi_X'(0)\big)^2.
		\]
	\end{proposition}
	
	\begin{remark}
		(Taylor expansion.)
		If \(\phi_X\) exists near \(0\), then formally
		\[
		\phi_X(t)=\mathbb{E}\big[e^{tX}\big]
		=\mathbb{E}\left[\sum_{n=0}^{\infty}\frac{t^nX^n}{n!}\right]
		=\sum_{n=0}^{\infty}\frac{t^n}{n!}\mathbb{E}[X^n],
		\]
		so m.g.f.\ is a generating function for moments.
Interchanging expectation and infinite sum needs justification; existence of m.g.f.\ near \(0\) is a standard sufficient condition.
	\end{remark}
	

	
	\begin{example}
		(Poisson.)
		If \(X\sim \mathrm{Poisson}(\lambda)\), then
		\begin{align*}
			\phi_X(t)
			&=\sum_{k=0}^{\infty} e^{tk}\frac{\lambda^k}{k!}e^{-\lambda}
			=e^{-\lambda}\sum_{k=0}^{\infty}\frac{(\lambda e^{t})^k}{k!}
			=e^{-\lambda}\exp(\lambda e^{t})
			=\exp\!\big(\lambda(e^{t}-1)\big).
		\end{align*}
		From this,
		\[
		\phi_X'(t)=\lambda e^t\exp\!\big(\lambda(e^t-1)\big),
		\qquad
		\phi_X'(0)=\lambda=\mathbb{E}X,
		\]
		and similarly \(\mathrm{Var}(X)=\lambda\).
	\end{example}
	
\begin{example}
	(Bernoulli and Binomial.)
	If \(X\sim \mathrm{Be}(p)\), then
	\[
	\phi_X(t)=\mathbb{E}[e^{tX}]=(1-p)e^{0}+pe^{t}=1-p+pe^{t}.
	\]
	If \(X\sim \mathrm{Bin}(n,p)\), then we may write
	\[
	X=\sum_{i=1}^n Y_i,
	\]
	where \(Y_1,\dots,Y_n\) are i.i.d.\ \(\mathrm{Be}(p)\). Hence
	\[
	\phi_X(t)
	=\mathbb{E}[e^{tX}]
	=\mathbb{E}\!\left[e^{t\sum_{i=1}^n Y_i}\right]
	=\mathbb{E}\!\left[\prod_{i=1}^n e^{tY_i}\right].
	\]
	By independence,
	\[
	\phi_X(t)
	=\prod_{i=1}^n \mathbb{E}[e^{tY_i}]
	=\prod_{i=1}^n \phi_{Y_i}(t).
	\]
	Since each \(Y_i\sim \mathrm{Be}(p)\),
	\[
	\phi_{Y_i}(t)=1-p+pe^t,
	\]
	and therefore
	\[
	\phi_X(t)=\bigl(1-p+pe^t\bigr)^n.
	\]
\end{example}
	

\begin{example}
	(Gaussian.)
	If \(X\sim N(\mu,\sigma^2)\), then
	\[
	\phi_X(t)=\mathbb{E}[e^{tX}]
	=\int_{-\infty}^{\infty} e^{tx}\frac{1}{\sqrt{2\pi}\sigma}
	\exp\!\left(-\frac{(x-\mu)^2}{2\sigma^2}\right)\,dx.
	\]
	Combining the exponentials and completing the square,
	\[
	tx-\frac{(x-\mu)^2}{2\sigma^2}
	=
	-\frac{\bigl(x-(\mu+\sigma^2 t)\bigr)^2}{2\sigma^2}
	+\mu t+\frac{\sigma^2 t^2}{2}.
	\]
	Therefore
	\[
	\phi_X(t)
	=
	\exp\!\left(\mu t+\frac{\sigma^2 t^2}{2}\right)
	\int_{-\infty}^{\infty}
	\frac{1}{\sqrt{2\pi}\sigma}
	\exp\!\left(
	-\frac{\bigl(x-(\mu+\sigma^2 t)\bigr)^2}{2\sigma^2}
	\right)\,dx.
	\]
	The integral equals \(1\), since it is the integral of a Gaussian density. Hence
	\[
	\phi_X(t)=\exp\!\left(\mu t+\frac{\sigma^2 t^2}{2}\right),
	\qquad t\in\mathbb{R}.
	\]
	
	Now let \(X_1,\dots,X_n\) be independent Gaussian random variables with
	\[
	X_i\sim N(\mu_i,\sigma_i^2),
	\qquad i=1,\dots,n.
	\]
	If
	\[
	S=\sum_{i=1}^n X_i,
	\]
	then by independence,
	\[
	\phi_S(t)
	=
	\prod_{i=1}^n \phi_{X_i}(t)
	=
	\prod_{i=1}^n
	\exp\!\left(\mu_i t+\frac{\sigma_i^2 t^2}{2}\right).
	\]
	So
	\[
	\phi_S(t)
	=
	\exp\!\left(
	t\sum_{i=1}^n \mu_i
	+\frac{t^2}{2}\sum_{i=1}^n \sigma_i^2
	\right).
	\]
	This is the MGF of
	\[
	N\!\left(\sum_{i=1}^n \mu_i,\ \sum_{i=1}^n \sigma_i^2\right).
	\]
	Therefore
	\[
	S\sim N\!\left(\sum_{i=1}^n \mu_i,\ \sum_{i=1}^n \sigma_i^2\right).
	\]
\end{example}


	
	\begin{proposition}[Application: sums of independent Poisson r.v.'s]
		If \(X\sim \mathrm{Poisson}(\lambda)\), \(Y\sim \mathrm{Poisson}(\mu)\), and \(X\perp\!\!\!\perp Y\), then
		\[
		X+Y\sim \mathrm{Poisson}(\lambda+\mu).
		\]
	\end{proposition}
	
	\begin{proof}
		Using the product rule,
		\[
		\phi_{X+Y}(t)=\phi_X(t)\phi_Y(t)
		=\exp(\lambda(e^t-1))\exp(\mu(e^t-1))
		=\exp((\lambda+\mu)(e^t-1)).
		\]
		This is exactly the m.g.f.\ of \(\mathrm{Poisson}(\lambda+\mu)\). By uniqueness of m.g.f.\ (on an interval around \(0\)),
		\(X+Y\sim \mathrm{Poisson}(\lambda+\mu)\).
	\end{proof}
	
	\begin{note}
	Next: since m.g.f.\ turns sums into products, it is also a key tool for limit theorems (LLN/CLT), where the number of summands tends to infinity.
	\end{note}
	
	
\section{Limit Theorems}

We have discussed sums of finitely many r.v.'s.
Now we send the number of summands to infinity and study the limiting behavior.



\begin{proposition}
	(Markov inequality.)
	For any non-negative r.v. \(X\) and any \(a>0\),
	\[
	\mathbb{P}(X\ge a)\le \frac{\mathbb{E}X}{a}.
	\]
\end{proposition}

\begin{proof}
	Since \(X\ge a\mathbf{1}_{\{X\ge a\}}\), taking expectation gives
	\[
	\mathbb{E}X \ge a\,\mathbb{E}\mathbf{1}_{\{X\ge a\}}
	= a\,\mathbb{P}(X\ge a).
	\]
	Therefore \(\mathbb{P}(X\ge a)\le \mathbb{E}X/a\).
\end{proof}

\begin{remark}
	(Consequence 1: apply Markov to \(|X|\).)
	For any r.v. \(X\) with \(\mathbb{E}|X|<\infty\) and any \(a>0\),
	\[
	\mathbb{P}(|X|\ge a)\le \frac{\mathbb{E}|X|}{a}.
	\]
\end{remark}

\begin{remark}
	(Consequence 2: apply Markov to \(g(X)\).)
	Let \(X\ge 0\) and let \(g:[0,\infty)\to[0,\infty)\) be increasing.
	Then for any \(a>0\),
	\[
	\mathbb{P}(X\ge a)=\mathbb{P}(g(X)\ge g(a))\le \frac{\mathbb{E}[g(X)]}{g(a)}.
	\]
	Typical choices: \(g(x)=x^2\), \(g(x)=e^{tx}\) for \(t>0\).
\end{remark}



\begin{proposition}
	(Chebyshev inequality.)
	Let \(X\) be any r.v. with finite mean \(\mu=\mathbb{E}X\) and finite variance \(\mathrm{Var}(X)\).
	Then for any \(a>0\),
	\[
	\mathbb{P}(|X-\mu|\ge a)\le \frac{\mathrm{Var}(X)}{a^2}.
	\]
\end{proposition}

\begin{proof}
	Apply Markov inequality to the non-negative r.v. \((X-\mu)^2\):
	\[
	\mathbb{P}(|X-\mu|\ge a)
	=\mathbb{P}\big((X-\mu)^2\ge a^2\big)
	\le \frac{\mathbb{E}[(X-\mu)^2]}{a^2}
	=\frac{\mathrm{Var}(X)}{a^2}.
	\]
\end{proof}



\begin{definition}
	(Convergence in probability.)
	A sequence of r.v.'s \((Z_n)_{n\ge 1}\) is said to converge to \(Z\) in probability if
	for every \(\varepsilon>0\),
	\[
	\lim_{n\to\infty}\mathbb{P}(|Z_n-Z|>\varepsilon)=0.
	\]
	We write \(Z_n \xrightarrow{\mathbb{P}} Z\).
\end{definition}



Let \(X_1,X_2,\dots\) be i.i.d.\ r.v.'s with common mean \(\mu\) and variance \(\sigma^2<\infty\).
Define the sample mean
\[
\overline{X}_n=\frac{1}{n}\sum_{i=1}^n X_i.
\]

\begin{theorem}
	(Weak Law of Large Numbers.)
	If \(X_1,X_2,\dots\) are i.i.d.\ with \(\mathbb{E}X_i=\mu\) and \(\mathrm{Var}(X_1)=\sigma^2<\infty\),
	then
	\[
	\overline{X}_n \xrightarrow{\mathbb{P}} \mu.
	\]
	Equivalently, for any fixed \(\varepsilon>0\),
	\[
	\lim_{n\to\infty}\mathbb{P}\!\left(\left|\frac{1}{n}\sum_{i=1}^n X_i-\mu\right|>\varepsilon\right)=0.
	\]
\end{theorem}

\begin{proof}
	First compute mean and variance.
	By linearity of expectation,
	\[
	\mathbb{E}\overline{X}_n
	=\frac{1}{n}\sum_{i=1}^n \mathbb{E}X_i
	=\mu.
	\]
	By independence and \(\mathrm{Var}(aX)=a^2\mathrm{Var}(X)\),
	\[
	\mathrm{Var}(\overline{X}_n)
	=\mathrm{Var}\!\left(\frac{1}{n}\sum_{i=1}^n X_i\right)
	=\frac{1}{n^2}\sum_{i=1}^n \mathrm{Var}(X_i)
	=\frac{1}{n^2}\cdot n\sigma^2
	=\frac{\sigma^2}{n}.
	\]
	Now apply Chebyshev inequality to \(\overline{X}_n\):
	\[
	\mathbb{P}(|\overline{X}_n-\mu|>\varepsilon)
	\le \frac{\mathrm{Var}(\overline{X}_n)}{\varepsilon^2}
	=\frac{\sigma^2}{n\varepsilon^2}\xrightarrow[n\to\infty]{}0.
	\]
\end{proof}

\begin{remark}
	Compare with the degenerate (completely dependent) case in which
	\[
	X_1=X_2=\cdots=X,
	\]
	where \(X\) is a random variable. Then
	\[
	\overline{X}_n=\frac{X_1+\cdots+X_n}{n}
	=\frac{X+\cdots+X}{n}
	=X
	\]
	for every \(n\). Thus the sample average is still just the same random variable \(X\), so in general it does not get closer to the constant \(\mathbb{E}[X]\) as \(n\to\infty\). In other words, no averaging occurs because the \(X_i\) contain no new independent information. This shows that the independence assumption is essential in the above proof.
\end{remark}


	Q: What is the limiting distribution of
	\[
	\frac{1}{\sqrt{n}}\sum_{i=1}^n (X_i-\mu)\,?
	\]
	This leads to the Central Limit Theorem.


\begin{theorem}
	(Central Limit Theorem.)
	Let \(X_1,X_2,\dots\) be i.i.d.\ with \(\mathbb{E}X_1=\mu\) and \(\mathrm{Var}(X_1)=\sigma^2\in(0,\infty)\).
	Define the normalized sum
	\[
	S_n:=\frac{1}{\sqrt{n}}\sum_{i=1}^n (X_i-\mu).
	\]
	Then, as \(n\to\infty\),
	\[
	S_n \xrightarrow{d} N(0,\sigma^2).
	\]
	Equivalently,
	\[
	\frac{\sqrt{n}\,(\overline{X}_n-\mu)}{\sigma}\xrightarrow{d} N(0,1).
	\]
\end{theorem}



\begin{proof}[Sketch of proof]
	Without loss of generality, assume \(\mu=0\) (apply CLT to \(X_i-\mu\)).
	Fix \(t\in\mathbb{R}\) and consider the m.g.f.\ of
	\[
	\frac{1}{\sqrt{n}}\sum_{i=1}^n X_i.
	\]
	By independence,
	\[
	\mathbb{E}\exp\!\left(t\cdot \frac{1}{\sqrt{n}}\sum_{i=1}^n X_i\right)
	=
	\prod_{i=1}^n \mathbb{E}\exp\!\left(\frac{t}{\sqrt{n}}X_i\right)
	=
	\left(\mathbb{E}e^{(t/\sqrt{n})X_1}\right)^n.
	\]
	Let \(\phi(u)=\mathbb{E}[e^{uX_1}]\) be the m.g.f.\ of \(X_1\) near \(u=0\).
	Using a Taylor expansion at \(u=0\) and \(\mathbb{E}X_1=0\), \(\mathbb{E}X_1^2=\sigma^2\),
	\[
	\phi(u)=1+\frac{\sigma^2}{2}u^2+o(u^2)
	\quad \text{as } u\to 0.
	\]
	Substitute \(u=t/\sqrt{n}\):
	\[
	\phi\!\left(\frac{t}{\sqrt{n}}\right)
	=
	1+\frac{\sigma^2}{2}\frac{t^2}{n}+o\!\left(\frac{1}{n}\right).
	\]
	Therefore,
	\[
	\left(\phi\!\left(\frac{t}{\sqrt{n}}\right)\right)^n
	=
	\left(1+\frac{\sigma^2 t^2}{2n}+o\!\left(\frac{1}{n}\right)\right)^n
	\longrightarrow
	\exp\!\left(\frac{\sigma^2 t^2}{2}\right).
	\]
	The limit \(\exp(\sigma^2 t^2/2)\) is the m.g.f.\ of \(N(0,\sigma^2)\).
	Hence the normalized sums converge in distribution to \(N(0,\sigma^2)\).
\end{proof}

	
	
\end{document}